% Appendix: attempted RWC witness on the Shrikhande / rook pair.

\section{An attempted RWC witness on the Shrikhande and $4 \times 4$
rook graphs}\label{app:rwc-witness}

This appendix supports Remark~\ref{rem:rwc} of Section~\ref{sec:measures} with
the numerical computation that motivated it. We attempted to
witness $\mathrm{RWC} \notin \mathcal{F}_{\wl}$ on the pair of
strongly regular graphs with parameters $(16, 6, 2, 2)$: the
Shrikhande graph $S$ and the $4 \times 4$ rook's graph $R$. The
attempt failed under every symmetric bipartition we tested. The
values below are exact absorbing-chain computations run through
the script at \texttt{experiments/src/rwc\_witness.py}.

\subsection{The two graphs}

The $4 \times 4$ rook's graph $R$ has vertex set $\{(i, j) : 0 \le
i, j \le 3\}$ and edges $\{\{(i, j), (i', j')\} : i = i' \text{ or }
j = j', (i, j) \neq (i', j')\}$. The Shrikhande graph $S$ has
vertex set $\Z/4\Z \times \Z/4\Z$ and edges $\{\{u, v\} : v - u \in
\{\pm (1, 0), \pm (0, 1), \pm (1, 1)\}\}$ under arithmetic modulo
$4$. Both graphs are $6$-regular; both have $16$ vertices and $48$
edges. The strongly regular parameters $(16, 6, 2, 2)$ are verified
by direct check in \texttt{rwc\_witness.py}. The two graphs are not
isomorphic: $R$ contains $K_4$ subgraphs (each row and each column)
while $S$ does not (rows in $S$ are $4$-cycles).

Under a constant node feature, both graphs collapse to a single
$1$-WL colour class of multiplicity $16$; the fingerprints agree
and $R \sim_{\wl} S$.

\subsection{The bipartition and featured $1$-WL check}

We use the bipartition $S_+ = \{(i, j) : i \in \{0, 1\}\}$, $S_- =
\{(i, j) : i \in \{2, 3\}\}$, so $|S_+| = |S_-| = 8$. Every vertex
of either graph has three same-side neighbours and three cross-side
neighbours in this bipartition. Encoded as a feature $x_v \in \{-1,
+1\}$, the depth-$1$ $1$-WL colour of a vertex is a function of
$(x_v, \text{same-side count}, \text{cross-side count})$. This
tuple is $(x_v, 3, 3)$ for every vertex in both graphs, and
refinement stabilises immediately. The featured fingerprints agree:
$(R, x) \sim_{\wl} (S, x)$.

Any MPNN with continuous update and readout gives $M(R, x) = M(S,
x)$ under this bipartition-as-feature, by Proposition~\ref{prop:upper}.

\subsection{Exact RWC computation}

We use $k = 1$ authority per side. Both graphs are vertex-transitive
so the specific choice within a side is immaterial; we take $(0, 0)$
as the $S_+$ authority and $(2, 2)$ as the $S_-$ authority in each
graph. A simple random walk from a uniform starting vertex of a
side terminates upon hitting either authority; $P_{st}$ is the
probability that a walk starting on side $s$ hits the authority of
side $t$ first.

The absorbing-chain computation solves $(I - Q) X = R$ where $Q$ is
the transition matrix among the $14$ non-absorbing vertices and $R$
records transitions into the two absorbing states. The solution
gives per-vertex absorption probabilities; averaging over starting
vertices of each side gives the four $P_{st}$ entries.

On both graphs the computation returns
\[
P_{++} = P_{--} = \frac{11}{21} \approx 0.5238,
\qquad
P_{+-} = P_{-+} = \frac{10}{21} \approx 0.4762,
\]
and hence
\[
\mathrm{RWC}(R) = \mathrm{RWC}(S) = P_{++} P_{--} - P_{+-} P_{-+} =
\frac{21}{441} = \frac{1}{21} \approx 0.0476.
\]
The two random walk controversy scores coincide to $10^{-4}$ in
float and exactly as rationals under symmetric bipartition and
authority choices.

\subsection{Why the witness fails on this pair}

Both graphs are vertex-transitive and cospectral (they share the
adjacency spectrum $\{6^1, 2^6, (-2)^9\}$ as SRGs with parameters
$(16, 6, 2, 2)$). Random walk absorption probabilities on a
vertex-transitive graph with a highly symmetric absorbing set are
determined by the walk's spectral decomposition; cospectral vertex-
transitive graphs with matching absorbing patterns produce matching
absorption probabilities. Under any bipartition that preserves the
symmetry (as the row-based bipartition does), the resulting RWC
values also coincide.

Constructing an $\mathrm{RWC}$ witness would require a pair of
$1$-WL equivalent featured graphs that are \emph{not} vertex-
transitive-symmetric under any suitable bipartition. The Cai-Fürer-
Immerman construction \citep{CFI1992} generates families of
$1$-WL equivalent non-isomorphic pairs, but the gadgets it uses
break vertex transitivity in ways that would still need to be
matched to a random-walk-controversy bipartition. We leave the
construction as an open technical problem; the taxonomy of
Section~\ref{sec:measures} therefore places $\mathrm{RWC}$ in a
"status open" row rather than in the not-$1$-WL-invariant column.
